% Notation for this section lives in gameops/notation.tex, which is shared
% with the TeXFrog HTML viewer so that the pseudocode renders identically in
% the paper and in the browser.

\subsection{A hint-aware Fiat--Shamir reduction}\label{sec:fs-hints}

The security statement of~\cite[Theorem~1]{C:ABDPW25} is obtained by composing
three results:
\begin{itemize}
    \item~\cite[Lemma~3.5]{EC:AABN02}, reduces the \EUFCMA security of $\SIG[\Sigma]$ to the \textsf{IMP-PA} security of $\Sigma$ at the price of a multiplicative factor $\numRO+1$; 
    \item~\cite[Lemma~3.1]{C:ABDPW25} then replaces the transcript oracle of the \textsf{IMP-PA} game by the hint-assisted simulator, at the price of an additive $\hinthvzkadv{\Sigma}{\hintdistr}{\Sim}{\cdot, \numOSign}$ term;
    \item~\cite[Lemma~3.2]{C:ABDPW25} that concludes by special soundness.
\end{itemize}
Because the second step is applied \emph{inside} the first one, the hint-assisted \textsf{wHVZK} advantage ends up being multiplied by $\numRO+1$.

This is an artefact of the modularisation rather than of the reduction: the transcript oracle of the \textsf{IMP-PA} game is only ever used to answer the signing queries of the forger, and there are $\numOSign$ of those, not $\numRO$.
In this section we prove a single lemma that performs the two steps in the opposite order, so that in the resulting bound $\hinthvzkadv{\Sigma}{\hintdistr}{\Sim}{\cdot, \numOSign}$ is a purely additive term.

\paragraph{Preliminaries.}
Throughout, $\Sigma = (\Prover_1, \Prover_2, \Verif)$ is a $\Sigma$-protocol for a relation $R$ with challenge space $\cC$, the family $\hintdistr = \set{\hintdistr_x}_x$ is a hint distribution for $R$ (\cite[Definition~3. 1]{C:ABDPW25}), and $\Sim$ is a $\hintdistr$-hint-assisted simulator for $\Sigma$ (\cite[Definition~3. 2]{C:ABDPW25}). 
We write $\SIG[\Sigma]$ for the Fiat--Shamir signature scheme obtained from $\Sigma$, in which the
challenge of a signature on $\msg$ under the public key $\pk = x$ is computed as
$\chall \gets \RO(\com \concat x \concat \msg)$.

For a game $\bG$ we write $\bG^\adv \Rightarrow 1$ for the event that the game outputs $1$ when run with $\adv$, and for $4 \leq i \leq 6$ we let $\goodev_i$ be the event that the $i$-th game never sets $\badev$.
Recall that two games are \emph{identical-until-$\badev$} if their code differs only in statements that follow the setting of the flag to \btrue, and recall the two standard consequences of this~\cite{EC:BelRog06}: 
\begin{itemize}
    \item for identical-until-$\badev$ games $\bG, \bG'$ and any adversary $\adv$, \begin{equation}\label{eq:fundamental} \abs{\Pr[\bG^\adv \Rightarrow 1] - \Pr[\bG'^\adv \Rightarrow 1]} \leq \Pr[\bG^\adv \text{ sets the flag}], \end{equation} where the right-hand side may equally be taken to be $\Pr[\bG'^\adv \text{ sets the flag}]$, since identical-until-$\badev$ games set the flag with the same probability; 
    \item and, denoting by $\goodev, \goodev'$ the events that the flag is never set, \begin{equation}\label{eq:good} \Pr[\bG^\adv \Rightarrow 1 \wedge \goodev] = \Pr[\bG'^\adv \Rightarrow 1 \wedge \goodev']. \end{equation}
\end{itemize}

Finally we write 
\begin{equation}
  \hintimppaadv{\Sigma}{\hintdistr}{\bdv}{\numOSign} \defeq
  \Pr\left[\hintdistr\text{-}\mathsf{hint}\text{-}\mathsf{IMP}\text{-}\mathsf{PA}_\Sigma(\bdv, \numOSign) = 1\right]
\end{equation}
for the advantage of a two-stage adversary $\bdv = (\bdv_1, \bdv_2)$ in the $\hintdistr$-hint-\textsf{IMP-PA} game of~\cite[Definition~3.5 and Game~3]{C:ABDPW25}.
Note that, unlike in the \textsf{IMP-PA} game, $\bdv$ has \emph{no} access to a transcript oracle, and that this game is not in the random oracle model: an adversary that needs a random oracle simulates it internally.

\paragraph{Statement.} 

\begin{lemma}[Hint-aware Fiat--Shamir]\label{lem:fs-hints}
    Let $(R, \hintdistr)$ be a relation with hints, let $\Sigma$ be a $\Sigma$-protocol for $R$ with challenge space $\cC$, and let $\Sim$ be a $\hintdistr$-hint-assisted simulator for $\Sigma$.
    For any PPT algorithm $\adv$ against the \EUFCMA security of $\SIG[\Sigma]$ making at most $\numRO$ queries to $\RO$ and at most $\numOSign$ queries to $\OSign$, there exist a two-stage PPT algorithm $\bdv = (\bdv_1, \bdv_2)$ and a PPT algorithm $\ddv$ such that
    \begin{multline} 
        \Advantage{\EUFCMA}{\SIG[\Sigma]}(\adv) \leq (\numRO + 1)\cdot \hintimppaadv{\Sigma}{\hintdistr}{\bdv}{\numOSign} + \\ 
        + \hinthvzkadv{\Sigma}{\hintdistr}{\Sim}{\ddv, \numOSign}  + (\numRO + \numOSign + 1)\,\numOSign \cdot \minentropy{\Sigma}.
    \end{multline}
    The running time of $\ddv$ is that of $\adv$ plus one evaluation of $\Ver_\Sigma$ and $O(\numRO + \numOSign)$ operations on the random-oracle table; 
    the running time of $\bdv$ is the same plus $\numOSign$ evaluations of $\Sim$.
\end{lemma}

\begin{proof}
    We use code-based game playing~\cite{EC:BelRog06}, following the proof of~\cite[Lemma~3.5]{EC:AABN02}; 
    the games are summarized in \Cref{fig:games}.
    \begin{figure}
        \centering 
        \tfrenderfigure{gameops}{G0,G1,G2,G3,G4,G5,G6}
        \caption{Games $\tfgamename{gameops}{G0}$--$\tfgamename{gameops}{G6}$ of the proof of Lemma~\ref{lem:fs-hints}.}
        \label{fig:games}
    \end{figure}

\subsubsection*{Normalising the forger.}
Exactly as in the proof of~\cite[Lemma~3.5]{EC:AABN02}, we replace $\adv$
by a forger $\adv'$ of essentially the same running time, making at most one
extra $\RO$ query, with
$\Advantage{\EUFCMA}{\SIG[\Sigma]}(\adv) \leq
\Advantage{\EUFCMA}{\SIG[\Sigma]}(\adv')$ and with the following three
properties, where $x$ denotes the public key that $\adv'$ receives:
\begin{enumerate}[label=(\roman*)]
    \item every $\RO$ query of $\adv'$ parses as $\com \concat x \concat
        \msg$, with middle component equal to $x$;
    \item before outputting a forgery $(\msg^*, \sig^* = (\com^*, \chall^*,
        \rsp^*))$, the algorithm $\adv'$ has queried $\RO$ on $\com^* \concat
        x \concat \msg^*$;
    \item $\msg^*$ was never queried to $\OSign$.
\end{enumerate}
% Indeed, $\adv'$ runs $\adv$ and forwards its queries, except that a query
% which does not parse as in (i) --- either because it is malformed or because
% its middle component differs from $x$ --- is answered from a separate table,
% sampled lazily and uniformly at random by $\adv'$ itself. Such points are
% never programmed by $\OSign$ and never enter the winning condition, so this
% does not change the view of $\adv$; using a table rather than fresh
% randomness keeps the answers consistent on repeated queries. Property (ii) is
% obtained by making the query $\RO(\com^* \concat x \concat \msg^*)$ just
% before halting, and property (iii) does not decrease the advantage, since a
% forgery on a signed message never wins the \EUFCMA game. The forger $\adv'$
% makes at most $\numRO + 1$ queries to $\RO$, all of them conforming; from now
% on we write $\adv$ for $\adv'$, and $\numRO$ still denotes the query bound of
% the original forger.

\subsubsection*{Game $\tfgamename{gameops}{G0}$.}
The \EUFCMA game of $\SIG[\Sigma]$, written so that the random oracle is
lazily sampled into a table $\HT$ and so that $\QT$ records, in order, the
arguments of the \emph{fresh} $\RO$ queries of $\adv$. Queries answered from
$\HT$ and reprogrammings performed by $\OSign$ are not recorded.

Lines $23$--$26$ set $\chall_\sctr \coloneqq \HT[u]$ if that entry is defined
and $\chall_\sctr \sample \cC$ otherwise; line $27$ then rewrites the same
value and is a no-op in the first case. This is exactly the signing algorithm
of $\SIG[\Sigma]$. Line $35$ reads $\HT[u^*]$ directly rather than calling
$\RO$, which is harmless because the normalised forger has always queried
that point.


\subsubsection*{Game $\tfgamename{gameops}{G1}$.}
Obtained from $\tfgamename{gameops}{G0}$ by deleting the single line $26$: a
signing query that hits an already defined table entry no longer re-uses the
stored value, but overwrites it with the freshly sampled challenge
$\chall_\sctr$.

Line $26$ follows the setting of $\badprog$ at line $25$, so the two games
are identical-until-$\badprog$. The commitment produced by $\Prover_1(x,w)$ is
independent of the message and of the history of the game, so the flag is set
with probability at most
$(\numRO + \numOSign + 1)\,\numOSign \cdot \minentropy{\Sigma}$.



\subsubsection*{Game $\tfgamename{gameops}{G2}$.}
The $\numOSign$ transcripts consumed by the signing oracle are sampled
\emph{eagerly}, in $\mathbf{Initialize}$, instead of on demand, and $\badprog$
is dropped. This is possible precisely because in
$\tfgamename{gameops}{G1}$ the value stored in $\HT$ no longer influences the
transcript, so $\pi_i$ is a function of $(x,w)$ and of fresh randomness only.

The rewriting is perfect. Its point is that $\mathbf{Initialize}$ now runs the
experiment $\Real_\Sigma(1^\secp, \numOSign)$ verbatim, which is what makes the
next hop a reduction to hint-assisted \textsf{wHVZK}.


\subsection*{Reduction $\tfgamename{gameops}{Red1}$ from $\tfgamename{gameops}{G2}$ to $\tfgamename{gameops}{G3}$.}
This reduction from \tfgamename{gameops}{G2} to \tfgamename{gameops}{G3} feeds the transcripts using the Hint-assisted wHVZK from \cite[Expderiment 2]{C:ABDPW25}.
Hence, by~\cite[Definition~3.2]{C:ABDPW25}, we have
\begin{equation}\label{eq:g2g3}
    \Pr[\tfgamename{gameops}{G2}^\adv \Rightarrow 1]
    \leq \Pr[\tfgamename{gameops}{G3}^\adv \Rightarrow 1]
    + \hinthvzkadv{\Sigma}{\hintdistr}{\Sim}{\ddv, \numOSign}.
\end{equation}
This is the step that differs from~\cite{C:ABDPW25}: the swap is performed on
the \EUFCMA game itself and no guessing has taken place yet, so the loss is a
single hint-assisted \textsf{wHVZK} advantage with $\numOSign$ hints.


\begin{figure}
    \centering 
    \tfrendergame[diff=G2]{gameops}{G3}
    \caption{Reduction $\ddv$ from 
        \tfgamename{gameops}{G2} to \tfgamename{gameops}{G3}, which uses the distinguisher $\ddv$ to turn the forger $\adv$ into a hint-assisted \textsf{wHVZK} distinguisher.
    }\label{fig:reduction}
\end{figure}


\subsubsection*{Game $\tfgamename{gameops}{G3}$.}
$\Real_\Sigma(1^\secp, \numOSign)$ is replaced by
$\HintSim_\Sigma(1^\secp, \numOSign, \Sim)$: the witness $w$ is discarded and
the $i$-th transcript is produced as $\pi_i \gets \Sim(x, \hint_i)$ with
$\hint_i \gets \hintdistr_x$.

This is the hop that differs from the reduction of Aardal et al. The swap is
performed on the \EUFCMA game itself and no guessing has taken place yet, so
the loss is a \emph{single} hint-assisted \textsf{wHVZK} advantage with
$\numOSign$ hints, rather than $\numRO + 1$ of them. The distinguisher needs
only the statement $x$ and the public $\Ver_\Sigma$: the witness is required
neither to answer the queries of $\adv$ nor to decide whether $\adv$ has won.

From here on the game does not use the witness at all.


\subsubsection*{Game $\tfgamename{gameops}{G4}$.}
Bookkeeping is added: an index $\fpguess \sample \set{1 \until \numRO+1}$ is drawn at line~\ref{gm:init-fp} and compared at line \ref{gm:fin-check-j} with the smallest index $j$ such that $\QT[j] = u^*$.

The index $\fpguess$ is drawn before both early rejections, hence on every run, and it is used neither in the answers to the queries of $\adv$ nor in the computation of the output.
Consequently the whole run of $\adv$ is independent of $\fpguess$, and hence so are both the event $\set{\tfgamename{gameops}{G4}^\adv \Rightarrow 1}$ and the index $j$.
In particular for any run of $\tfgamename{gameops}{G4}$ that outputs $1$ any of the $\numRO +1$ shifted combination of $j,\fpguess$ is equally likely.
By line~\ref{gm:fin-check-j} $\goodev_4$ holds if and only if $j = \fpguess$.
Putting all this together gives
\begin{equation}\label{eq:g3g4}
    \begin{aligned}
        \Pr[\tfgamename{gameops}{G3}^\adv \Rightarrow 1]
            &= \Pr[\tfgamename{gameops}{G4}^\adv \Rightarrow 1]
            = (\numRO+1)\cdot
            \Pr[\tfgamename{gameops}{G4}^\adv \Rightarrow 1 \wedge j = \fpguess] 
            \\
            &= (\numRO+1)\cdot
            \Pr[\tfgamename{gameops}{G4}^\adv \Rightarrow 1 \wedge \goodev_4].
    \end{aligned}
\end{equation}

    % Note that \eqref{eq:g3g4} must be phrased through the joint event, and
    % \emph{not} in the factorised form $\Pr[\bG \Rightarrow 1 \wedge \goodev] =
    % \Pr[\bG \Rightarrow 1]\cdot\Pr[\goodev]$ used in~\cite{EC:AABN02}: because of
    % the early rejections at lines \ref{gm:fin-check-msg} and \ref{gm:fin-check-chal}, which have no counterpart there,
    % the flag $\badev$ is also left unset on runs that reject early, so
    % $\Pr[\goodev_4] > 1/(\numRO+1)$.


\subsubsection*{Game $\tfgamename{gameops}{G5}$.}
Additionally overwrites $\chall^*$ with $\HT[\QT[\fpguess]]$ when $j \neq \fpguess$, that is, it runs the verifier on the challenge sitting at the \emph{guessed} position rather than at the correct one.
This modification (line \ref{gm:fin-overwrite}) follows the setting of $\badev$ at line \ref{gm:fin-set-bad}, so $\tfgamename{gameops}{G4}$ and $\tfgamename{gameops}{G5}$ are identical-until-$\badev$:
\begin{equation}\label{eq:g4g5}
    \Pr[\tfgamename{gameops}{G4}^\adv \Rightarrow 1 \wedge \goodev_4]
    = \Pr[\tfgamename{gameops}{G5}^\adv \Rightarrow 1 \wedge \goodev_5].
\end{equation}


\subsubsection*{Game $\tfgamename{gameops}{G6}$.}
Game $\tfgamename{gameops}{G6}$ finally draws the challenge $\chall' \sample \cC$ before calling the adversary $\adv$ in line~\ref{gm:adv-run}. 
Also, $\chall'$ is given as input as the answer to the $\fpguess$-th fresh $\RO$ query. 
This is a change of sampling order only.
In $\tfgamename{gameops}{G5}$ the answers to the fresh $\RO$ queries are independent uniform elements of $\cC$ and $\fpguess$ is an independent uniform index; in $\tfgamename{gameops}{G6}$ for every fixed value of $\fpguess$, the answers to the fresh $\RO$ queries are again independent uniform elements of $\cC$. 
Thus
\begin{equation}\label{eq:g5g6}
    \Pr[\tfgamename{gameops}{G5}^\adv \Rightarrow 1 \wedge \goodev_5]
    = \Pr[\tfgamename{gameops}{G6}^\adv \Rightarrow 1 \wedge \goodev_6].
\end{equation}


\subsubsection*{The impersonator.}
Game $\tfgamename{gameops}{G6}$ is exactly the simulation carried out by the following two-stage adversary $\bdv = (\bdv_1, \bdv_2)$ against the $\hintdistr$-hint-\textsf{IMP-PA} security of $\Sigma$.
On input a statement $x$ and hints $\hint_1 \until \hint_\numOSign \gets \hintdistr_x$, the algorithm $\bdv_1$ computes $\pi_i \gets \Sim(x, \hint_i)$ for $i = 1 \until \numOSign$, draws $\fpguess \sample \set{1 \until \numRO+1}$, and runs $\adv$ on $\pk = x$, answering its queries as in $\tfgamename{gameops}{G6}$, until $\adv$ makes its $\fpguess$-th fresh $\RO$ query. 

It parses that query as $\com^* \concat x \concat \msg^*$ and outputs $\com^*$ together with the state of the simulation. 
On input the challenge $\chall^\circ \sample \cC$, the algorithm $\bdv_2$ sets $\HT[\com^* \concat x \concat \msg^*] \coloneqq \chall^\circ$, returns $\chall^\circ$ to $\adv$ as the answer to its $\fpguess$-th fresh query, resumes the simulation, and outputs the response $\rsp^*$ contained in the forgery of $\adv$.
Since $\chall^\circ$ is drawn uniformly at random from $\cC$, the view of $\adv$ is independent, as long as it is returned after the $\fpguess$-th fresh query.

It is clear that on the event $\tfgamename{gameops}{G6}^\adv \Rightarrow 1 \wedge \goodev_6$ the algorithm $\bdv_2$ can return $\resp^*$ such that the triple $(\com^*, \chall^\circ, \rsp^*)$, is a valid transcript, that is the winning condition of $\bdv$. 
Consequently
\begin{equation}\label{eq:g6b}
    \Pr[\tfgamename{gameops}{G6}^\adv \Rightarrow 1 \wedge \goodev_6]
    \leq \hintimppaadv{\Sigma}{\hintdistr}{\bdv}{\numOSign}.
\end{equation}

Chaining all the game hops gives the claim.
\end{proof}


\subsection{The improved \EUFCMA bound}\label{sec:eufcma-gen}

Combining Lemma~\ref{lem:fs-hints} with the special-soundness step
of~\cite[Lemma~3.2]{C:ABDPW25} we obtain the announced generalisation
of~\cite[Theorem~1]{C:ABDPW25}.

\begin{theorem}\label{thm:EUFCMA-gen}
    Let $(R, \hintdistr)$ be a relation with hints.
    Let $\Sigma$ be a $\Sigma$-protocol for the relation $R$ that has challenge
    space $\cC$ and is special-sound with respect to a soundness relation
    $\tilde{R}$.
    Additionally, let $\Sim$ be a $\hintdistr$-hint-assisted simulator for
    $\Sigma$.
    For any PPT algorithm $\adv$ against the \EUFCMA security of $\SIG[\Sigma]$
    there exist a PPT algorithm $\ddv$ and an expected polynomial time algorithm
    $\cE$ such that
    \begin{align*}
        \Advantage{\EUFCMA}{\SIG[\Sigma]}(\adv)
        &\leq \left(\numRO+1\right)\cdot
        \left(\hintreladv{\tilde{R}}{\hintdistr}{\cE}{\numOSign}
        + \frac{1}{\lvert \cC \rvert}\right) \\
        &\phantom{\leq ~}
        + \hinthvzkadv{\Sigma}{\hintdistr}{\Sim}{\ddv, \numOSign}
        + (\numRO + \numOSign + 1)\,\numOSign \cdot \minentropy{\Sigma},
    \end{align*}
    where $\numRO$ and $\numOSign$ are upper bounds on the number of queries that
    $\adv$ makes to $\RO$ and $\OSign$, respectively.
\end{theorem}

\begin{proof}
    Let $\bdv$ and $\ddv$ be the algorithms given by Lemma~\ref{lem:fs-hints}, so
    that
    \[
        \begin{aligned}
            \Advantage{\EUFCMA}{\SIG[\Sigma]}(\adv)
            &\leq (\numRO+1)\cdot\hintimppaadv{\Sigma}{\hintdistr}{\bdv}{\numOSign}
            + \hinthvzkadv{\Sigma}{\hintdistr}{\Sim}{\ddv, \numOSign} \\
            &\phantom{\leq ~}
            + (\numRO+\numOSign+1)\numOSign\cdot\minentropy{\Sigma}.
        \end{aligned}
    \]
    Since $\Sigma$ is special sound with respect to $\tilde{R}$,
    \cite[Lemma~3.2]{C:ABDPW25} applied to the two-stage adversary $\bdv$
    produces an expected polynomial time algorithm $\cE$ for the relation with
    hints $(\tilde{R}, \hintdistr)$ such that
    \[
        \hintimppaadv{\Sigma}{\hintdistr}{\bdv}{\numOSign}
        \leq \hintreladv{\tilde{R}}{\hintdistr}{\cE}{\numOSign}
        + \frac{1}{\abs{\cC}}.
    \]
    Substituting gives the claim. If a strict PPT algorithm is required, one
    applies~\cite[Corollary~3.1]{C:ABDPW25} to $\cE$: this yields a PPT algorithm
    $\cE'$ with $\hintreladv{\tilde{R}}{\hintdistr}{\cE}{\numOSign} \leq 2 \cdot
    \hintreladv{\tilde{R}}{\hintdistr}{\cE'}{\numOSign}$, so that the first term
    of the bound becomes $(\numRO+1)\bigl(2 \cdot
    \hintreladv{\tilde{R}}{\hintdistr}{\cE'}{\numOSign} + 1/\abs{\cC}\bigr)$.
    Note that $\cE'$ runs in polynomial time in $\secp$ only when the success
    probability of $\cE$ is the inverse of a polynomial. \qed
\end{proof}

\begin{remark}[Comparison with~{\cite[Theorem~1]{C:ABDPW25}}]\label{rem:comparison}
    The bound of~\cite[Theorem~1]{C:ABDPW25} reads
    \[
        \begin{aligned}
            \Advantage{\EUFCMA}{\SIG[\Sigma]}(\adv)
            &\leq (\numRO+1)\cdot
            \left(\hintreladv{\tilde{R}}{\hintdistr}{\cE}{\numOSign}
            + \hinthvzkadv{\Sigma}{\hintdistr}{\Sim}{\bdv, \numOSign}
            + \frac{1}{\abs{\cC}}\right) \\
            &\phantom{\leq ~}
            + (\numRO+\numOSign+1)\numOSign\cdot\minentropy{\Sigma}.
        \end{aligned}
    \]
    Theorem~\ref{thm:EUFCMA-gen} improves on it by taking the hint-assisted
    \textsf{wHVZK} term out of the factor $\numRO+1$; all the other terms, the
    number of hints and the running times of the constructed algorithms are
    unchanged, so the improvement comes for free. It matters whenever the
    hint-assisted simulator is only \emph{statistically} close to the honest
    prover with a gap that is not overwhelmingly small. This is exactly the
    situation for the odd-degree variant of \sqisign{}: the hint distribution has
    a failure branch, taken with probability at most $1/\epsilon$, on which no
    odd-degree response isogeny of degree below $\responseupperbound$ exists, and
    this branch contributes a term $\numOSign/\epsilon$ to
    $\hinthvzkadv{\Sigma}{\hintdistr}{\Sim}{\ddv, \numOSign}$ on top of the terms
    that are already present in the even-degree case. In the bound
    of~\cite[Theorem~1]{C:ABDPW25} that contribution would instead be
    $(\numRO+1)\,\numOSign/\epsilon$. For the same security claim, $\log\epsilon$
    may therefore be smaller by $\log \numRO$ bits, which with the NIST bound
    $\numRO \leq 2^{64}$ is what makes the smaller values of $\log \epsilon$ in
    Table~\ref{tab:timing} admissible.
\end{remark}

\begin{remark}[Why this order of the game hops]\label{rem:order}
    The reordering performed in the proof of Lemma~\ref{lem:fs-hints} is
    constrained from both sides, and it is not free.
    \begin{enumerate}[label=(\alph*)]
        \item The replacement of the honest transcripts by simulated ones must
            come \emph{after} the hop $\tfgamename{gameops}{G0} \to \tfgamename{gameops}{G1}$. In
            $\tfgamename{gameops}{G0}$ the challenge of a signature is the value of the
            random oracle at $\com \concat x \concat \msg$, which need not be
            uniform and independent of the rest, since $\adv$ may have queried
            that point before. The experiments $\Real_\Sigma$ and
            $\HintSim_\Sigma$ of~\cite[Experiment~2]{C:ABDPW25} produce
            transcripts with \emph{independent uniform} challenges, so they can
            only be plugged in once the signing oracle has been made to sample
            the challenge itself and to reprogram the random oracle accordingly.
            This is what the min-entropy term pays for.
        \item The replacement must come \emph{before} the guessing hops
            $\tfgamename{gameops}{G3} \to \tfgamename{gameops}{G6}$, since those are what introduce
            the factor $\numRO+1$.
    \end{enumerate}
    In particular, Lemma~\ref{lem:fs-hints} cannot be obtained by composing
    \cite[Lemma~2.2]{C:ABDPW25} with \cite[Lemma~3.1]{C:ABDPW25} as black boxes:
    the game hop has to be inserted \emph{inside} the proof
    of~\cite[Lemma~3.5]{EC:AABN02}, between the games that isolate the
    reprogramming failure and the games that perform the guess. As a by-product,
    the min-entropy of the \emph{simulated} commitment distribution is never
    needed: from $\tfgamename{gameops}{G1}$ on the signing oracle reprograms the random
    oracle unconditionally, so only $\minentropy{\Sigma}$, the min-entropy of the
    honest protocol, appears in the bound.
\end{remark}
